% Generated draft: requires mathematical and source review.
\documentclass[11pt]{article}
\usepackage[margin=25mm]{geometry}
\usepackage{fontspec}
\setmainfont{Latin Modern Roman}
\usepackage{amsmath,amssymb,mathtools,unicode-math}
\setmathfont{Latin Modern Math}
\usepackage{enumitem,needspace,xurl,hyperref}
\setlength{\parindent}{0pt}
\setlength{\parskip}{5pt}
\setlength{\emergencystretch}{3em}
\newcommand{\N}{\mathbb{N}}
\newcommand{\Z}{\mathbb{Z}}
\newcommand{\Q}{\mathbb{Q}}
\newcommand{\R}{\mathbb{R}}
\newcommand{\C}{\mathbb{C}}
\newcommand{\F}{\mathbb{F}}
\newcommand{\A}{\mathbb{A}}
\newcommand{\PP}{\mathbb{P}}
\newcommand{\E}{\mathbb{E}}
\newcommand{\eps}{\varepsilon}
\newcommand{\dd}{\,\mathrm{d}}
\begin{document}
\section*{UnsolvedMath ID: 30003074}
\textbf{Source problem code:} catalogue entry
\subsection{Definitions and mathematical statement}
The supplied record asks whether a particular Lefschetz property is determined by the combinatorics of an arrangement. It does not specify the base field, its characteristic, the ambient polynomial ring, or the admissible values of the exponent. The following is a conditional formulation using the usual maximal-rank interpretation of the Lefschetz terminology; these missing parameters require source verification.

Fix an ambient standard graded polynomial ring \(R=k[x_0,\ldots,x_n]\), with \(\deg x_i=1\). Here \(k\) and \(n\) are placeholders for the intended ambient data, not additional universally quantified parameters of a verified conjecture. For a homogeneous ideal \(K\subseteq R\), write \(A_K=R/K\) and let \((A_K)_d\) denote its degree-\(d\) component. Under the usual interpretation, having the Strong Lefschetz Property at range \(2\) in degree \(b-2\) means that, for a general linear form \(L\in R_1\), the multiplication map
\[
\mu_{K,L}\colon (A_K)_{b-2}\longrightarrow (A_K)_b,
\qquad \overline f\longmapsto\overline{L^2f},
\]
has maximal rank, namely
\[
\operatorname{rank}(\mu_{K,L})
=\min\{\dim_k(A_K)_{b-2},\dim_k(A_K)_b\}.
\]
Over an infinite field, the existence of such an \(L\) is equivalent to this condition holding on a nonempty Zariski-open subset of \(R_1\). The record does not specify the field or its convention for the word ``general.'' This displayed formulation directly applies when \(b\geq2\); the intended lower bound on \(b\) is not supplied.

For nonzero linear forms \(\ell_1,\ldots,\ell_{2b+1}\), a standard interpretation of their arrangement is the central hyperplane arrangement
\[
\mathcal A_{\ell}=\{\ker(\ell_i)\subseteq k^{n+1}:1\leq i\leq2b+1\}.
\]
Its intersection lattice consists of all intersections of subfamilies, including the ambient space as the intersection of the empty subfamily, ordered by reverse inclusion. Under the usual interpretation, two arrangements have the same combinatorics when their intersection lattices are isomorphic. The source record does not explicitly define this convention or identify the central or projective setting.

With these interpretations and the intended ambient parameters fixed, the question is whether, for every pair of lists of linear forms \((\ell_i)_{i=1}^{2b+1}\) and \((h_i)_{i=1}^{2b+1}\), the following implication holds. Set
\[
I=(\ell_1^b,\ldots,\ell_{2b+1}^b),
\qquad
J=(h_1^b,\ldots,h_{2b+1}^b).
\]
Assume that both displayed generating lists are minimal, meaning that none of their generators belongs to the ideal generated by the others, and that \(\mathcal A_{\ell}\) and \(\mathcal A_h\) have isomorphic intersection lattices. Must maximal rank of \(\mu_{I,L}\) for general \(L\) imply maximal rank of \(\mu_{J,M}\) for general \(M\in R_1\)? The general forms \(L\) and \(M\) need not be identical.

The supplied background concerns line arrangements, suggesting a projective-plane setting, but the statement itself does not specify it. No assertion for arbitrary dimensions or characteristics, and no additional Artinian hypothesis on the quotients, is supplied or established here.
\subsection{Short English statement}
If an ideal minimally generated by \(2b+1\) powers of linear forms has the specified degree-wise Lefschetz property, must every such ideal with the same arrangement combinatorics have it too? The intended field, dimension, exponent range, and precise conventions need verification. This is a statement-only draft requiring human mathematical review.
\subsection{Sources}
\begin{itemize}
\item Supplied local UnsolvedMath registry record: permanent ID \(30003074\), title ``Combinatorial Invariance of Strong Lefschetz Properties,'' exact source problem code \texttt{catalogue entry}, in the Oberwolfach Reports Open Problems collection. The record gives the proposed year as 2016 and identifies no proposer.
\item Supplied source citation: \emph{Mini-Workshop: Arrangements of Subvarieties, and their Applications in Algebraic Geometry} (2016), Oberwolfach Reports, pp.~653--702, DOI \texttt{10.4171/owr/2016/14}. Source URL and supplied literature link ``OWR report 14/2016'': \url{https://doi.org/10.4171/owr/2016/14}. The background describes these reports as records of mathematical workshops and their open-problem discussions; no more precise location within the report is supplied.
\item Additional supplied literature link, labeled ``Pagaria'': \url{https://doi.org/10.1112/jlms.12829}. The record supplies no title, full author information, or pinpoint reference for this item.
\item Supplied literature assessment dated exactly \texttt{2026-08-21}: the record labels the problem open and states that this conjecture is explicitly equivalent to Terao's conjecture for line arrangements, that current literature still treats that conjecture as open, and that many subclasses are settled. This is supplied dated curation, not an independently verified assessment. The background explicitly calls for independent verification of time-sensitive status and cited claims.
\item The supplied statement status is \texttt{corrected\_verified}, with the explanation that notation, typography, and sentence structure were repaired without changing mathematical content. This registry label does not establish independent mathematical verification of the formulation presented here.
\end{itemize}
\end{document}
